\section{Conclusion}
\label{sec:conclusion}

This paper develops a stochastic structural framework for valuing Bitcoin treasury vehicles---publicly traded companies whose primary asset is BTC, acquired through a leverage flywheel of equity, debt, and preferred stock issuance. We jointly model BTC price dynamics, the equity-over-NAV premium as an Ornstein--Uhlenbeck process, BTC holding accumulation as a compound Poisson process linked to the flywheel, and share-count evolution from issuance and conversion.

From this framework we derive five diagnostic indicators: the Implied BTC Growth Rate (IBGR), Implied Leverage Elasticity (ILE), Total Equity Elasticity (TEE), Premium Mean-Reversion Index (PMRI), and Implied Funding Requirement Distribution (IFRD). Together, these indicators decompose the valuation of a BTC treasury vehicle into its structural components---balance-sheet leverage, sentiment premium, accretion efficiency, and financing sustainability---providing investors and analysts with a rigorous toolkit for assessing these novel entities.

Calibrating to Strategy (MSTR) through December 2025, we find that the flywheel has been highly accretive (18.3\% annualized BTC growth, nearly identical per-share), balance-sheet leverage is currently moderate (ILE = 1.17), and the capital structure supports a three-year survival probability exceeding 99.4\%. The premium currently sits 1.5 standard deviations below its long-run mean, suggesting potential for mean-reversion.

These results are conditioned on the calibration sample and the model's assumptions of parameter stationarity. As Strategy's capital structure continues to evolve---with multiple preferred stock series, expanding ATM programs, and potential international vehicles---the framework presented here provides a foundation for increasingly granular analysis. The extension to multi-layer preferred capital stacks, outlined in the discussion, is a natural and important direction for future work.
